\title{Direct Preference Optimization: Your Language Model is Secretly a Reward Model}

\begin{abstract}
Although large-scale unsupervised language models (LMs) acquire extensive world knowledge and demonstrate some reasoning capabilities, exerting fine-grained control over their behavior remains challenging because of their entirely unsupervised training paradigm. Current approaches to enhance this controllability involve gathering human assessments of the comparative quality among model outputs and then fine-tuning the LMs so that they better reflect these expressed preferences, typically employing reinforcement learning from human feedback (RLHF). Nonetheless, RLHF is a multi-stage and frequently unstable approach that first constructs a reward model mirroring human judgements, then uses reinforcement learning to adapt the base LM for maximizing the estimated reward while attempting to remain close to the original model’s distribution. In this work, we propose a novel way of parameterizing the reward model in RLHF that allows us to analytically derive the optimal policy, thereby reducing the conventional RLHF problem to the minimization of a standard classification loss. Our method, termed \textit{Direct Preference Optimization} (DPO), offers a reliable, efficient, and resource-friendly alternative, obviating the need for LM sampling or extensive hyperparameter search during fine-tuning. Experimental results reveal that DPO is capable of aligning LMs to human preferences as effectively, if not better, than current state-of-the-art techniques. In particular, DPO surpasses PPO-based RLHF methods in sentiment control tasks, and provides competitive or enhanced performance in summarization and single-turn dialogue benchmarks, all while offering greater simplicity in both implementation and training.
\end{abstract}

\section{Introduction}
Large unsupervised language models (LMs) trained on massive corpora have demonstrated unexpected competencies~\citep{chowdhery2022palm, brown2020language, touvron2023llama, bubeck2023sparks}. Despite this, these models are developed on data reflecting a vast spectrum of human objectives, priorities, and expertise. Not all of these human-derived attributes are suitable for replication; for instance, it is desirable for an AI coding assistant to \textit{recognize} typical programming errors to offer corrections, yet when generating code, the model should be predisposed towards the exceptional, albeit less frequent, instances of expert coding found in its data. Similarly, while it is beneficial for a language model to \textit{know about} a widespread misconception held by half of people, we do not want it to present this mistaken belief as accurate in half of relevant interactions. Therefore, the ability to choose the model's \emph{preferred outputs and conduct} from the breadth of its \textit{knowledge and skills} is vital for ensuring that AI systems remain safe, effective, and amenable to control \citep{ouyang2022training}. Standard approaches tend to align LMs with human preferences by leveraging reinforcement learning (RL); however, we demonstrate that the same RL-driven objective employed by current techniques can be solved precisely using a straightforward binary cross-entropy loss, significantly streamlining the process of learning preferences.

In broad terms, prevailing techniques instill specific behavioral traits in language models via carefully constructed datasets that capture human preferences regarding desired conduct. This preference alignment is implemented subsequent to the initial phase of extensive unsupervised pre-training on large-scale textual data. While supervised fine-tuning on human-generated examples of high-quality responses is the most direct method for preference learning, the most effective methodologies employ reinforcement learning from human (or AI-generated) feedback (RLHF/RLAIF; \citep{christiano2017deep,bai2022constitutional}). RLHF strategies fit a reward function to preference-labeled data and apply RL to adjust the language model’s policy, encouraging responses that score highly under the learned reward model while constraining deviation from the original model. Although RLHF yields models with strong conversational and coding capabilities, it introduces considerable complexity relative to supervised learning—requiring the training of multiple models and repeated policy sampling during optimization, all of which escalate computational overhead.

In this work, we present a method for optimizing language models to better match human preferences, doing so without the need for explicit reward modeling or reinforcement learning approaches. We introduce \textit{{Direct Preference Optimization} (DPO)}, an algorithm that inherently seeks the same reward maximization objective as traditional RLHF methods, including a KL-divergence penalty, but offers a much more accessible implementation and training process. The core mechanism of DPO is to boost the relative log-likelihood of preferred responses compared to non-preferred alternatives, while employing an adaptive, example-specific importance weighting. This adjustment mitigates the model collapse issues encountered with naïve objectives based solely on probability ratios. As with prior approaches, DPO is grounded in a formal preference model (e.g., the Bradley-Terry model; \cite{bradley1952rankanalysis}) for assessing the alignment between reward functions and observed preference judgments. However, in contrast to methods that first use this preference model to learn a reward function and then optimize a policy against that learned function, DPO reparameterizes the preference loss so that it is defined directly in terms of the policy parameters. With access to a dataset consisting of human comparative preferences among model outputs, DPO can train a policy via a straightforward binary cross entropy loss, arriving at an optimal policy for an implicitly defined reward that fits the preference data.

The primary contribution of this work is Direct Preference Optimization (DPO), a lightweight algorithm that bypasses reinforcement learning in favor of direct, preference-based fine-tuning of language models. Our empirical results indicate that DPO matches or surpasses the effectiveness of prior approaches—including PPO-based RLHF—across various preference-driven applications, such as sentiment control, summarization, and dialogue generation, and is effective even with large-scale models up to 6B parameters.

\section{Related Work}

Self-supervised language models of greater scale demonstrate the capacity to handle certain tasks without prior examples (zero-shot) \citep{radford2019language}, or with minimal context via prompt-based learning (few-shot) \citep{gpt3,megatron,chowdhery2022palm}. Nonetheless, their effectiveness on downstream applications and alignment with user instructions can be substantially enhanced through fine-tuning with corpora comprising instructions and human-generated outputs \citep{mishra-etal-2022-cross,sanh2022multitask,chung2022scaling,thoppilan2022lamda}. This approach, known as `instruction-tuning,' permits large language models (LLMs) to extrapolate to previously unseen instructions, thereby broadly improving model utility \citep{chung2022scaling}. Although instruction tuning has demonstrated substantial benefits, acquiring \textit{relative} evaluations from humans (comparing responses) is often simpler than curating expert-level demonstrations. Consequently, subsequent studies have employed human preference datasets to further fine-tune LLMs, resulting in improved performance across tasks such as translation \citep{kreutzer-etal-2018-reliability}, summarization \citep{stiennon2022learning,ziegler2020finetuning}, story generation \citep{ziegler2020finetuning}, and instruction adherence \citep{ouyang2022training,ramamurthy2023is}. Typically, these methods begin by training a neural reward model to reflect the dataset of preferences, commonly within a framework like the Bradley-Terry preference model \citep{bradley1952rankanalysis}. Afterward, reinforcement learning algorithms—such as REINFORCE \citep{williams1992reinforce}, proximal policy optimization (PPO; \cite{schulman2017proximal}), or related techniques \citep{ramamurthy2023is}—are used to adjust the language model in order to maximize the computed reward. Another closely associated research trajectory utilizes instruction-tuned LLMs with human feedback for the automatic generation of supplementary synthetic preference data, targeting specific attributes like safety or benign behavior \citep{bai2022constitutional}. This is accomplished using only loosely-supervised input from humans via textual guidelines to inform the LLM's assessments. Collectively, these strategies reflect the convergence of two research streams: reinforcement learning approaches for language model training on various objectives~\citep{Ranzato2015SequenceLT,paulus2018a,wu2018learning}, and broader frameworks for preference-based learning from humans \citep{christiano2017deep,kupcsik2018learning}. Despite the practical appeal of leveraging comparative human judgments, implementing reinforcement learning-based fine-tuning on large language models presents significant hurdles; this paper introduces a theoretically grounded method for optimizing relative preferences that circumvents the need for reinforcement learning.

The task of inferring policies from preferences, independent of language applications, has been investigated in both bandit frameworks and reinforcement learning, leading to a variety of methodological proposals. In particular, contextual bandit learning using comparative judgments—where actions are ranked or compared instead of assigned explicit reward values—is referred to as contextual dueling bandits (CDB; \cite{yue2012karmed,dudik2015contextual}). Within this paradigm, where absolute reward feedback is missing, the theoretical analysis replaces the concept of an optimal policy with the \textit{von Neumann winner}, which describes a policy expected to outperform \textit{every} other policy at least half the time \citep{dudik2015contextual}. A key distinction in CDBs is that feedback is delivered sequentially (online), whereas human preference-based learning often utilizes a fixed, pre-collected set of action pairs annotated with offline preferences \citep{yan2022human}. Likewise, \textit{preference-based RL} (PbRL) algorithms rely on binary preference feedback produced by an \textit{unknown} scoring function instead of explicit rewards \citep{BusaFekete2014,ruiz2023dueling}. Multiple approaches to PbRL have been advanced, with some enabling efficient reuse of off-policy data; however, these typically involve a two-step process of first inferring a latent reward or scoring function and subsequently optimizing it \citep{jain2013learning,BusaFekete2014,christiano2017deep,sadigh2017active,kupcsik2018learning}. In contrast, we introduce a direct policy learning method that forgoes explicit reward estimation in favor of optimizing a policy to align with observed preferences in a single phase.

\section{Preliminaries}\label{section:prelims}

We summarize the RLHF process as outlined by \citeauthor{ziegler2020finetuning}, with extensions in later works (\citep{stiennon2022learning, bai2022training, ouyang2022training}). The RLHF framework generally comprises three primary steps: (1) supervised fine-tuning (SFT), (2) preference elicitation and reward modeling, and (3) reinforcement learning optimization.

\textbf{SFT}: The RLHF pipeline is typically initiated by performing supervised learning-based fine-tuning on a pre-trained language model using task-specific high-quality datasets (such as for dialogue or summarization), producing the model $\pi^\text{SFT}$.

\textbf{Reward Modelling Phase}: In the following stage, the SFT-tuned model is provided with prompts $x$, generating answer pairs $(y_1, y_2)\sim \pi^\text{SFT}(y \mid x)$. These responses are then evaluated by human annotators, who indicate a preference—represented as $y_w\succ y_l \mid x$, where $y_w$ is the favored answer over $y_l$ for the input $x$. The actual decision process behind these preferences is assumed to depend on a latent, unobserved reward model $r^*(y, x)$. Several frameworks exist for modeling such preferences, with the Bradley-Terry (BT) model \cite{bradley1952rankanalysis} being a frequently adopted method (although the approach is extensible to other models such as the Plackett-Luce ranking models \citep{plackett1975analysis, luce2012individual} in cases with more than two ranked choices). According to the BT formulation, the distribution of human preferences $p^*$ can be described as follows:

\begin{equation}\label{eq:bradley-terry}
    p^*(y_1\succ y_2 \mid x)=\frac{\exp\left(r^*(x, y_1)\right)}{\exp\left(r^*(x, y_1)\right) + \exp\left(r^*(x, y_2)\right)}.
\end{equation}

Given a fixed dataset of pairwise comparisons, $\mathcal{D}=\bigl\{x^{(i)}, y_w^{(i)}, y_l^{(i)}\bigr\}_{i=1}^N$, generated according to $p^*$, one can model the reward function $r_{\phi}(x, y)$ parametrically and fit its parameters by maximizing the likelihood. By viewing the task as a binary classification, the associated negative log-likelihood is defined as follows:

\begin{equation}\label{eq:reward_model}
    \mathcal{L}_R(r_{\phi}, \mathcal{D}) = -\mathbb{E}_{(x, y_w, y_l)\sim \mathcal{D}}\bigl[\log \sigma(r_{\phi}(x, y_w)- r_{\phi}(x, y_l))\bigr]
\end{equation}

with $\sigma$ denoting the logistic sigmoid. For language models, the architecture for $r_{\phi}(x, y)$ typically leverages weights initialized from the SFT policy $\pi^\text{SFT}(y \mid x)$ and appends a linear projection to the last transformer block, producing a single-valued reward estimate \cite{ziegler2020finetuning}. To reduce variability in the reward signal, previous approaches standardize the reward so that  $\mathbb{E}_{x,y\sim \mathcal{D}}\left[r_\phi(x, y)\right] = 0$ for any input $x$.

\textbf{RL Fine-Tuning Phase}: At this stage, the reward function is used as feedback for training the language model. Consistent with prior literature~\citep{jaques2017sequence, jaques2020human}, the training is cast as the following optimization:

\begin{equation}\label{eq:RL}
\max_{\pi_{\theta}}  \mathbb{E}_{x\sim \mathcal{D}, y\sim \pi_{\theta}(y \mid x)}\bigl[r_{\phi}(x, y)\bigr] - \beta\mathbb{D}_{\textrm{KL}}\bigl[\pi_{\theta}(y\mid x)\mid \mid \pi_\text{ref}(y\mid x)\bigr],
\end{equation}

where $\beta$ determines the extent to which the optimized policy $\pi_{\theta}$ diverges from a fixed reference $\pi_\text{ref}$, which is generally chosen as the SFT model $\pi^\text{SFT}$. Practically, $\pi_\theta$ is commonly initialized from $\pi^\text{SFT}$ as well. The regularization term serves to keep the updated policy near the distribution where the reward model is reliable, preserving response variety and guarding against collapsing onto a small set of maximally rewarded outputs. Since the language generation process is inherently discrete, this objective does not allow direct gradient computation and is usually tackled via reinforcement learning algorithms. Typically, the reward is rewritten as ${r(x, y) = r_{\phi}(x, y) -\beta (\log \pi_{\theta}(y\mid x) - \log \pi_\text{ref}(y\mid x))}$ and optimized using PPO \cite{schulman2017proximal}, as outlined in previous work \citep{ziegler2020finetuning, stiennon2022learning, bai2022training, ouyang2022training}.

\section{Direct Preference Optimization}\label{sec:DPO}

Recognizing the practical difficulties of scaling reinforcement learning techniques to large models, especially during language model fine-tuning, we propose a streamlined approach for optimizing policies via preference data. Departing from the traditional RLHF paradigm, which requires first inferring a reward model and then training a policy via RL, our method takes advantage of a strategic reward function parameterization that admits a closed-form expression for the optimal policy—eliminating the necessity of an RL training cycle. The crux of our method is an explicit analytical relationship that translates reward functions to their corresponding optimal policies, thereby enabling us to rephrase a loss on rewards as a loss on policies. This transformation circumvents the need to learn a separate reward function, yet still optimizes with respect to established models of preference data, such as the Bradley-Terry formulation

\textbf{Reformulating the DPO Objective.} We begin by considering the standard RL objective as outlined in previous research, specifically Eq.~\ref{eq:RL}, for an arbitrary reward function $r$. As established in works such as~\citep{peters2007reinforcement, peng2019advantage, korbak2022reinforcement, go2023aligning}, the optimal solution to the reward maximization problem with a KL divergence constraint, given in Eq.~\ref{eq:RL}, is:

\begin{equation}\label{eq:op_policy}
    \pi_r(y\mid x) = \frac{1}{Z(x)}\pi_\text{ref}(y\mid x)\exp\left(\frac{1}{\beta}r(x, y)\right),
\end{equation}

where the normalization factor $Z(x) =\sum_{y}\pi_\text{ref}(y\mid x)\exp\left(\frac{1}{\beta}r(x, y)\right)$ is known as the partition function (see Appendix \ref{app:derivation1} for a detailed derivation). In practice, even if we estimate $r$ by a maximum likelihood estimator $r_{\phi}$ of the true reward $r^*$, computing $Z(x)$ remains computationally burdensome~\citep{korbak2022reinforcement, go2023aligning}. This limitation reduces the tractability of directly working with this expression. Instead, we can manipulate Eq.~\ref{eq:op_policy} to write the reward in terms of its optimal policy $\pi_r$, the reference policy $\pi_\text{ref}$, and $Z(x)$: by taking logarithms and rearranging, we obtain

\begin{equation}\label{eq:main_eq}
    r(x,y) =\beta \log \frac{\pi_r(y\mid x)}{\pi_\text{ref}(y\mid x)} + \beta \log Z(x).
\end{equation}

This transformation applies both to the ground-truth reward $r^*$ and its associated optimal policy $\pi^*$. Importantly, the Bradley-Terry model considers only reward differences between two responses; that is, $p^*(y_1 \succ y_2 \mid x) = \sigma(r^*(x, y_1) - r^*(x, y_2))$. When we substitute the above parameterization from Eq.~\ref{eq:main_eq} into the preference model Eq.~\ref{eq:bradley-terry}, the partition function $Z(x)$ drops out, allowing us to represent the preference probability solely via the optimal policy $\pi^*$ and the reference policy $\pi_\text{ref}$. Consequently, the optimal RLHF policy $\pi^*$ within the Bradley-Terry framework fulfills the relation:

\begin{equation}\label{eq:objective}
    p^*(y_1\succ y_2 \mid x)=\frac{1}{1 + \exp\left(\beta \log \frac{\pi^*(y_2\mid x)}{\pi_\text{ref}(y_2\mid x)} - \beta \log \frac{\pi^*(y_1\mid x)}{\pi_\text{ref}(y_1\mid x)}\right)}
\end{equation}

The step-by-step derivation for this result can be found in Appendix~\ref{app:derivation2}. Although Eq.~\ref{eq:objective} is based on the Bradley-Terry framework, similar formulations may be deduced for the more general Plackett-Luce models~\citep{plackett1975analysis, luce2012individual}, as elaborated in Appendix~\ref{app:plackett_luce_models}.

Since we have now established a preference likelihood expressed directly in terms of the optimal policy (rather than the reward model), it becomes possible to define a maximum likelihood training objective over policy parameters $\pi_\theta$. In a manner analogous to Eq.~\ref{eq:reward_model} (the reward modeling objective), the policy objective can thus be formulated as:

\begin{equation}\label{eq:optimum_model}
    \mathcal{L}_\text{DPO}(\pi_{\theta}; \pi_\text{ref}) = -\mathbb{E}_{(x, y_w, y_l)\sim \mathcal{D}}\left[\log \sigma \left(\beta \log \frac{\pi_{\theta}(y_w\mid x)}{\pi_\text{ref}(y_w\mid x)} - \beta \log \frac{\pi_{\theta}(y_l\mid x

In this formulation, we infer an implicit reward by adopting a different parameterization, with its optimal policy represented directly by $\pi_\theta$. Additionally, since our approach mirrors the reparameterized Bradley-Terry model, it inherits several theoretical guarantees, such as consistency given appropriate assumptions regarding the preference data distribution \cite{bong2022generalized}. We provide a more detailed discussion of DPO’s theoretical underpinnings in Section~\ref{sec:theory}, particularly in comparison to existing literature.

\textbf{Interpreting the DPO update.} To build intuition for DPO, it is instructive to examine the gradient of the DPO objective, $\mathcal{L}_\text{DPO}$. The gradient with respect to model parameters $\theta$ is expressed as:

\begin{multline*}\label{eq:gradient}
    \nabla_\theta \mathcal{L}_\text{DPO}(\pi_\theta;\pi_\text{ref}) = \\ -\beta\mathbb{E}_{(x, y_w, y_l) \sim \mathcal{D}} \bigg[\underbrace{\sigma(\hat{r}_\theta(x, y_l) - \hat{r}_\theta (x, y_w))}_\text{larger influence when reward prediction is poor}\bigg[\underbrace{\nabla_\theta\log \pi(y_w \mid x)}_\text{boost preference for $y_w$} - \underbrace{\nabla_\theta\log\pi(y_l \mid x)}_\text{lower preference for $y_l$}\bigg]\bigg],
\end{multline*}

Here, $\hat{r}_\theta(x, y) = \beta \log \frac{\pi_\theta(y \mid x)}{\pi_\text{ref}(y \mid x)}$ denotes the reward implicitly encoded by $\pi_\theta$ relative to the reference $\pi_\text{ref}$ (refer to Section~\ref{sec:theory} for more detail). In essence, the loss gradient encourages the model to increase the probability of generating favored responses $y_w$ while reducing that of less preferred responses $y_l$. Notably, each sample’s influence is scaled according to how much the implicit reward model $\hat{r}_\theta$ inappropriately scores the disfavored completion over the preferred one, modulated by $\beta$. This weighting reflects how severely the model misranks the completions while incorporating the KL regularization’s impact. Our empirical findings highlight the necessity of this weighting term, as omitting it can lead to model collapse, as detailed in Appendix Table~\ref{tab:unlikelihood_generations}.

\textbf{Overview of DPO.} The DPO framework proceeds through the following steps: (1) For each prompt $x$, generate two candidate completions $y_1, y_2 \sim \pi_\text{ref}(\cdot \mid x)$ and annotate them according to human preferences, constructing the offline preference dataset $\mathcal{D} = \{x^{(i)}, y_w^{(i)}, y_l^{(i)}\}_{i=1}^N$; (2) update the language model $\pi_\theta$ by minimizing $\mathcal{L}_\text{DPO}$, with respect to $\pi_\text{ref}$, $\mathcal{D}$, and the chosen $\beta$. In applied scenarios, practitioners often prefer leveraging publicly available preference datasets instead of collecting fresh samples and human labels. Because these datasets are typically generated using $\pi^\text{SFT}$, we set $\pi_\text{ref} = \pi^\text{SFT}$ when possible. If $\pi^\text{SFT}$ is not present, we instead estimate $\pi_\text{ref}$ by maximizing the likelihood over preferred completions ${(x, y_w)}$, that is, by choosing ${\pi_\text{ref} = \argmax_{\pi}\mathbb{E}_{x, y_w \sim \mathcal{D}}\left[\log \pi(y_w \mid x)\right]}$. This approach helps reduce the gap between the actual (unobserved) reference distribution and the $\pi_\text{ref}$ implemented by DPO. Additional information on implementation strategies and hyperparameter choices is available in Appendix~\ref{app:implementation}.

\section{Theoretical Analysis of DPO} In the following, we deepen the conceptual understanding of DPO, offering theoretical justification, and contrast the strengths of DPO against limitations found in actor-critic approaches used in RLHF—such as PPO~\cite{schulman2017proximal}.

\label{sec:theory}

\subsection{Your Language Model Is Secretly a Reward Model} DPO eliminates the necessity of both fitting an explicit reward function and conducting RL for policy learning by employing a unified maximum likelihood training objective. Observe that the optimization in Eq.~\ref{eq:main_eq} is identical to a Bradley-Terry model where rewards are parameterized as $r^*(x, y) = \beta \log\frac{\pi^*_\theta(y \mid x)}{\pi_\text{ref}(y \mid x)}$, and the parameterization of $\pi_\theta$ directly mirrors the structure of reward model optimization given in Eq.~\ref{eq:reward_model}, after a suitable reparametrization. In what follows, we formalize the theoretical basis for this parameterization, demonstrate that it imposes no restriction on the set of attainable reward models, and prove that this approach enables perfect recovery of the optimal policy. We start by introducing an equivalence relation among reward functions.

\begin{definition}
Two reward functions $r(x, y)$ and $r'(x, y)$ are considered equivalent if and only if there exists a function $f$ such that ${r(x, y)-r'(x, y) = f(x)}$.     
\end{definition}

This definition readily establishes an equivalence relation, thereby partitioning the set of all reward functions into equivalence classes. This observation gives rise to the following lemmas:

\begin{lemma}\label{lemma:same_prefrence} Within the Plackett-Luce framework—including the Bradley-Terry model—reward functions belonging to the same equivalence class result in identical preference distributions.
\end{lemma}

\begin{lemma}\label{lemma:same_policy}
    Within the context of constrained RL, any two reward functions from a shared equivalence class generate the same optimal policy.
\end{lemma}

The demonstrations of these lemmas are elementary and appear in Appendix \ref{app:lemma1}. The initial lemma reflects a classical indeterminacy inherent in the Plackett-Luce models \cite{plackett1975analysis}. Owing to this identifiability limitation, it is typically necessary to impose further constraints to derive guarantees regarding the maximum likelihood estimates in Eq. \ref{eq:reward_model} \cite{bong2022generalized}. The second lemma asserts that every reward function within an equivalence class yields an identical optimal policy; thus, for optimization purposes, our objective reduces to recovering any representative reward from the optimal class. The following theorem, whose proof is found in Appendix~\ref{app:thm1}, formalizes this insight:

\begin{theorem}\label{thm:main}
    Provided certain mild assumptions, any reward equivalence class consistent with the Plackett-Luce (and, specifically, the Bradley-Terry) models admits a representation of the form ${r(x, y) = \beta \log \frac{\pi(y\mid x)}{\pi_\text{ref}(y\mid x)}}$ for some candidate model $\pi(y\mid x)$ and some fixed reference model $\pi_\text{ref}(y \mid x)$.
\end{theorem}

\begin{sproof}
    Take an arbitrary reward function $r(x, y)$, which defines an associated optimal policy $\pi_r(y \mid x)$ as given in Eq. \ref{eq:op_policy}. We now demonstrate that it is possible to construct an equivalent reward function from the equivalence class of $r$ in the reparameterized form indicated. Define the projection mapping $f$ by  
\begin{equation}
    f(r; \pi_\text{ref}, \beta)(x, y) = r(x, y) - \beta\log\sum_{y}\pi_\text{ref}(y\mid x)\exp\left(\frac{1}{\beta}r(x, y)\right)
\end{equation}
Here, $f$ acts by centering $r$ using the log-partition function (which is solely dependent on $x$) with respect to $\pi_r$. As the adjustment term depends only on $x$, $f(r; \pi_\text{ref}, \beta)(x, y)$ remains in the same equivalence class as $r(x, y)$. Substituting the expression for $r$ as per Eq.~\ref{eq:main_eq} (a general result), one obtains $f(r; \pi_\text{ref}, \beta)(x, y) = \beta \log \frac{\pi_r(y\mid x)}{\pi_\text{ref}(y\mid x)}$. Therefore, the projection $f$ returns an equivalent reward of the prescribed structure, confirming that this reparameterization encompasses all reward classes without any loss of generality.
\end{sproof}

Theorem~\ref{thm:main} can alternatively be interpreted as identifying the precise reward function selected by the DPO reparameterization within each equivalence class. Specifically, it selects the reward function that fulfills the following criterion:

\begin{equation}\label{eq:lag_p}
     \sum_{y}\underbrace{\pi_\text{ref}(y\mid x)\exp\left(\frac{1}{\beta}r(x, y)\right)}_{=\pi(y\mid x)\text{, using Thm.~\ref{thm:main} reparam.}} = 1,
\end{equation}

meaning that $\pi(y\mid x)$ constitutes a proper probability distribution, as it remains nonnegative and normalizes to one. Furthermore, building on Eq.~\ref{eq:op_policy}, it becomes evident that Eq.~\ref{eq:lag_p} corresponds to the partition function for the optimal policy associated with the reward function $r(x, y)$. The central innovation of the DPO framework is its ability to enforce particular constraints on the otherwise under-specified Plackett-Luce (and, in particular, Bradley-Terry) models of preference, thereby preserving the set of admissible reward functions while ensuring that the optimal policy in Eq.~\ref{eq:op_policy} is analytically tractable for every prompt $x$.

\subsection{Instability of Actor-Critic Algorithms} Our framework can further be leveraged to elucidate sources of instability in conventional actor-critic methods applied within RLHF pipelines, such as PPO. Here, we concentrate on the RL fine-tuning phase described in Section \ref{section:prelims}, aligning our perspective with the control-as-inference framework~\cite{levine2018reinforcement} relevant to the constrained RL problem introduced in Eq.~\ref{eq:RL}. We posit a parameterized policy $\pi_{\theta}(y\mid x)$ and seek to minimize $\mathbb{D}_{\text{KL}}[\pi_{\theta}(y|x) \mid\mid \pi^*(y\mid x)]$, where $\pi^*$ is the optimal policy (cf. Eq.~\ref{eq:optimum_model}) generated by the reward function $r_{\phi}(y, x)$. A straightforward manipulation yields the following optimization problem:

\begin{equation}\label{eq:AC}
    \max_{\pi_{\theta}}\mathbb{E}_{\pi_{\theta}(y\mid x)}\bigg[\underbrace{r_{\phi}(x, y) -\beta\log\sum_{y}\pi_\text{ref}(y\mid x)\exp\left(\frac{1}{\beta}r_{\phi}(x, y)\right)}_{f(r_{\phi}, \pi_\text{ref}, \beta)} - \underbrace{\beta\log\frac{\pi_{\theta}(y\mid x)}{\pi_\text{ref}(y\mid x)}}_{\text{KL}}\bigg]
\end{equation}

The objective being optimized here matches that of previous studies 
\citep{ziegler2020finetuning, stiennon2022learning, bai2022training, ouyang2022training}, where a DPO-equivalent reward corresponding to the $r_{\phi}$ reward family is employed. Within this framework, the normalization component in $f(r_{\phi}, \pi_\text{ref}, \beta)$ can be interpreted as representing the soft value function for the reference policy $\pi_\text{ref}$. Although this normalization term does not alter the optimal policy, omitting it can result in a policy gradient with significant variance, thereby compromising learning stability. One approach to address the normalization is to incorporate a trainable value function, though this can present its own optimization challenges. Previous methods have alternatively employed reward normalization by leveraging human completion baselines, which essentially provide a single-sample Monte Carlo approximation for the normalization factor. In contrast, the DPO reparameterization defines a reward function inherently independent of baseline estimations.

\section{Experiments}
Here, we conduct empirical studies to assess DPO’s effectiveness in learning policies directly from preference data. Initially, we utilize a controlled text-generation scenario to investigate how DPO balances the dual objectives of reward maximization and KL-divergence minimization relative to the reference policy. We compare this to standard preference optimization approaches, including PPO. Subsequently, DPO's scalability is examined in the context of more complex and larger-scale RLHF applications, such as tasks in summarization and dialogue. Our findings indicate that, with minimal hyperparameter adjustment, DPO consistently matches or outperforms established baselines—such as RLHF with PPO and the selection of the best among $N$ samples according to a learned reward model. Preceding our discussion of results, we provide a detailed outline of the experimental procedures, with further specifics included in Appendix~\ref{app:exp_details}.

\textbf{Tasks.} Our experimental evaluation encompasses three distinct open-ended text generation tasks. In each scenario, the employed algorithms optimize a policy based on a preference dataset $\mathcal{D}=\bigl\{x^{(i)}, y_w^{(i)}, y_l^{(i)}\bigr\}_{i=1}^N$. In the case of \textbf{controlled sentiment generation}, $x$ denotes the initial segment of a movie review from the IMDb dataset \cite{maas-EtAl:2011:ACL-HLT2011}, and the objective is for the policy to generate a continuation $y$ exhibiting positive sentiment. For the purpose of a controlled assessment, preference pairs for generations are produced using a sentiment classifier that has been pre-trained, ensuring that $p(\text{positive}\mid x,y_w)>p(\text{positive}\mid x,y_l)$. For SFT, we continue GPT-2-large's training until it converges, using samples from the training split of the IMDB dataset (see App~\ref{app:sentiment_details} for additional information). 

In the \textbf{summarization} task, $x$ represents a Reddit forum post, and the policy is tasked with generating a summary $y$ distilling the main ideas of the original post. Consistent with earlier research, we adopt the Reddit TL;DR summarization dataset \citep{volske-etal-2017-tl} in combination with a collection of human-labeled preferences gathered by \citeauthor{stiennon2022learning}. For this purpose, we utilize an SFT model that has been fine-tuned on summaries of forum posts written by humans\footnote{\url{https://huggingface.co/CarperAI/openai_summarize_tldr_sft}}, with RLHF training performed through the TRLX \citep{leandro_von_werra_2023_7790115} framework. The human preferences dataset originates from work by \citeauthor{stiennon2022learning} using samples from an SFT model trained in a similar fashion, though not identically. 

The third scenario, \textbf{single-turn dialogue}, frames $x$ as a human-issued prompt, which may encompass a broad spectrum of queries, from technical scientific inquiries to requests for advice. Here, the policy is required to craft a response $y$ that is both informative and engaging. We make use of the Anthropic Helpful and Harmless dialogue dataset \citep{bai2022training}, which comprises 170k dialogues conducted between humans and an automated assistant. Every dialogue culminates with two responses generated by a large, albeit unspecified, language model and a label specifying which response humans favored. Given the absence of a ready-made SFT model in this context, we construct the SFT model by fine-tuning an existing language model solely on the human-preferred completions.

\textbf{Evaluation.} We assess our algorithms through two distinct evaluation strategies. To scrutinize each algorithm's capability in solving the constrained reward maximization problem, we conduct controlled sentiment generation experiments where we compute the performance frontier in terms of achieved reward versus KL-divergence from the reference policy; here, direct computation is possible due to access to an explicit ground-truth reward (via a sentiment classifier). However, because the true reward function remains unavailable in real-world settings, we instead report each algorithm’s \textit{win rate} relative to a baseline policy. This assessment leverages GPT-4 as a stand-in for human evaluators, rating the quality of summaries and the helpfulness of dialogue responses in the summarization and single-turn dialogue tasks, respectively. Specifically, test-set reference summaries act as baselines for summarization, while for dialogue, the human-preferred responses from the test set serve this role. Despite recent evidence that LMs can surpass traditional metrics in evaluation \citep{Chen2023ExploringTU}, we validate our GPT-4-based protocol via a human annotation study described in Sec.~\ref{sec:human-judgments}, showing that GPT-4 judgments exhibit strong alignment with those from humans, often matching or exceeding the consistency among human annotators themselves.

\textbf{Methods.} Beyond DPO, we benchmark a variety of established methods for aligning language models with human preferences. As straightforward baselines, we investigate zero-shot prompting with \textbf{GPT-J} \citep{gpt-j} for summarization and 2-shot prompting with \textbf{Pythia-2.8B} \citep{biderman2023pythia} for dialogue. Our analysis also includes the \textbf{SFT} model and \textbf{Preferred-FT}, which is created by supervised fine-tuning on the selected completion $y_w$—drawn either from SFT outputs (for controlled sentiment and summarization) or a generic language model (for single-turn dialogue). Another approach is \textbf{Unlikelihood}~\citep{welleck2019neural}, which encourages the model to increase the probability of $y_w$ and reduce that of $y_l$, modulated by a tunable `unlikelihood' weight $\alpha \in [0, 1]$. We further evaluate \textbf{PPO} \citep{schulman2017proximal} using a learned reward function based on preference data, and its oracle counterpart \textbf{PPO-GT}, which relies on the true reward in the controlled sentiment scenario. For sentiment experiments, we employ two PPO-GT variants: an out-of-the-box implementation \cite{leandro_von_werra_2023_7790115} and a custom version with reward normalization and optimized hyperparameters, applying these enhancements also to PPO with learned rewards. Finally, we include the \textbf{Best of $N$} baseline, where $N$ responses are sampled from SFT (or Preferred-FT for dialogue), with the top candidate selected using a reward model trained on preferences. Though this strategy can yield strong results by decoupling reward model quality from PPO optimization, it proves computationally costly since it requires generating $N$ completions per input at inference.

\subsection{How well can DPO optimize the RLHF objective?}

In standard RLHF frameworks, the KL-regularized reward maximization objective is designed to maximize the reward while simultaneously ensuring that the policy does not diverge significantly from the reference policy. As a consequence, any comparison of methods must consider both the achieved reward and the degree of KL divergence; surpassing previous rewards at the expense of substantially higher KL may not be favorable. Figure~\ref{fig:frontier-tldr-main} illustrates the trade-off between reward and KL for different approaches in the sentiment classification scenario. For each algorithm, we conduct several experiments, systematically varying the policy conservativeness parameter in each trial (target KL $\in\{3,6,9,12\}$ for PPO, $\beta \in \{0.05,0.1,1,5\}$, $\alpha\in\{0.05,0.1,0.5,1\}$ for unlikelihood, random seeds for preferred-FT), culminating in a total of 22 runs. At intervals of 100 training steps up to convergence, we assess every policy on a fixed set of test prompts, measuring both the mean reward—using the actual reward function—and the mean sequence-level KL\footnote{Specifically, the total of the KL divergences across all timesteps.} relative to the reference policy, denoted as $\text{KL}\left(\pi\mid \mid \pi_\text{ref}\right)$. Our experiments reveal that DPO achieves a markedly superior frontier: it attains the highest rewards while maintaining low KL values. This is significant for several reasons. Firstly, DPO and PPO target the same objective, yet DPO is considerably more effective—its reward/KL balance uniformly exceeds that of PPO. Secondly, DPO still surpasses PPO in this trade-off curve, \emph{even under conditions where PPO has access to the true reward values} (PPO-GT).

\subsection{Can DPO scale to real preference datasets?}
\label{sec:dpo-real-datasets}
We proceed to assess the fine-tuning capabilities of DPO in tasks related to summarization and single-turn dialogue. In summarization, automatic metrics such as ROUGE have demonstrated weak alignment with human preferences~\citep{stiennon2022learning}. Previous research indicates that optimizing LMs using PPO with human preference data leads to more satisfactory summaries. For our evaluation, we generate completions on the test partition of the TL;DR summarization dataset and calculate each method's average win rate in comparison to the reference completions. Completions are drawn for all models across a range of sampling temperatures from 0.0 to 1.0, with results summarized in Figure~\ref{fig:frontier-tldr-main} (right). DPO, PPO, and Preferred-FT are all fine-tuned starting from the identical GPT-J SFT model\footnote{\url{https://huggingface.co/CarperAI/openai_summarize_tldr_sft}}. Experimental findings reveal that DPO achieves a win rate near 61\% at a temperature of 0.0, surpassing PPO, which reaches about 57\% at its optimal temperature. DPO also establishes a higher peak win rate relative to the best-of-$N$ baseline. It should be highlighted that DPO's $\beta$ parameter was not subject to significant tuning, implying that these findings may not fully represent DPO's capabilities. Additionally, DPO demonstrates considerably greater resilience to temperature variations compared to PPO, whose effectiveness declines to match the original GPT-J model at elevated temperatures. In contrast, Preferred-FT does not yield a notable improvement over the base SFT model. Finally, a direct comparison of DPO and PPO through human judgment, discussed in Section~\ref{sec:human-judgments}, shows that samples generated by DPO at temperature 0.25 are favored 58\% of the time over those produced by PPO at temperature 0.

For single-turn dialogue evaluation, we assess each method using a subset of the Anthropic HH dataset's test split \citep{bai2022training} containing only single-step interactions between humans and assistants. In GPT-4-based assessments, win rates are calculated by referencing the preferred completions from the test data to benchmark each approach. Due to the absence of a widely accepted SFT model for this particular task, our workflow begins with a pre-trained Pythia-2.8B. We apply Preferred-FT to create a reference model using the preferred completions, ensuring the completions are within the distribution of the model. Subsequently, we fine-tune the model with DPO. Additional comparisons are made with the Best of 128 Preferred-FT completions—the performance of the Best of $N$ baseline saturates at 128 completions for this dataset, as illustrated in Appendix Figure~\ref{fig:best-of-n}—and with a 2-shot prompted configuration of the Pythia-2.8B base model. DPO achieves equal or superior results, matching the best-performing temperature settings for each baseline. An RLHF model, trained via PPO on the Anthropic HH dataset \footnote{\url{https://huggingface.co/reciprocate/ppo_hh_pythia-6B}} and sourced from a reputable repository \footnote{\url{https://github.com/CarperAI/trlx/tree/main/examples/hh}}, was also tested; however, it could not surpass the performance of the Pythia-2.8B base model, regardless of the prompting strategy or temperature. Our prior findings from the TL;DR benchmark, combined with the observation that both PPO and Best of 128 optimize the same reward function, suggest that Best of 128 is a reasonable stand-in for PPO performance. In summary, among the tested approaches, DPO uniquely stands out as a computationally efficient method that not only surpasses the preferred completions in the Anthropic HH dataset, but also performs on par with, or better than, the more resource-intensive Best of 128 baseline. Furthermore, Figure~\ref{fig:dialogue-main} demonstrates that DPO rapidly achieves peak performance.

\subsection{Generalization to a new input distribution}

To further examine how PPO and DPO cope with distribution shifts, we assess the policies trained in our Reddit TL;DR summarization experiment on a new domain—specifically, news articles from the test split of the CNN/DailyMail dataset \citep{nallapati-etal-2016-abstractive}. For this evaluation, we retain the optimal sampling temperatures found on TL;DR (0 and 0.25). The comparative results are shown in Table~\ref{tab:ood}. To determine performance, we calculated GPT-4’s win rate by comparing model-generated summaries to the ground-truth ones, applying the same GPT-4 (C) prompt as before, except substituting ``forum post'' with ``news article.'' Notably, DPO policies sustain a clear advantage over PPO in this out-of-distribution context. This finding suggests that DPO policies have generalization capabilities on par with those of PPO, despite DPO not relying on the extra unlabeled Reddit TL;DR prompts leveraged by PPO.

\subsection{Validating GPT-4 judgments with human judgments}
\label{sec:human-judgments}
We conduct a human evaluation study to test the consistency of GPT-4’s assessments by leveraging outcomes from the TL;DR summarization task and two variants of GPT-4 prompting. The \textbf{GPT-4 (S)} (simple) prompt queries which summary best captures the key information from the post. Meanwhile, the \textbf{GPT-4 (C)} (concise) prompt extends this by also soliciting which summary is more concise; this prompt is tested because, under the \textbf{GPT-4 (S)} prompt, GPT-4 exhibits a preference for summaries that are lengthier and more repetitive compared to human raters. Appendix~\ref{app:prompts} details the prompts. We select three systems for comparison, representing the highest-performing (DPO, temp. 0.25), lowest-performing (PPO, temp. 1.0), and an intermediate (SFT, temp. 0.25) model, all matched against PPO’s best-performing greedy sampling baseline; this selection covers a spectrum of summary qualities. Our analysis reveals that both GPT-4 prompts yield human agreement rates comparable to human-human inter-annotator agreement, suggesting that GPT-4’s evaluations can serve as a reasonable stand-in for human judgment (due to the limited number of human raters, multiple annotations are only collected for the DPO and PPO-1 conditions). Overall, the \textbf{GPT-4 (C)} prompt delivers win rates most aligned with human preference; therefore, it is chosen as the main evaluation in Section~\ref{sec:dpo-real-datasets}. Further methodological specifics, including rater interface design and the human participants list, can be found in Appendix~\ref{app:human-study}.

\section{Discussion}
Learning from preferences constitutes an effective and scalable methodology for developing language models that are both powerful and aligned with user intent. In this work, we presented DPO, a straightforward approach for preference-based language model training that does not rely on reinforcement learning. Instead of forcing preference learning into a traditional RL framework to utilize existing RL solutions, DPO establishes a connection between the policy of a language model and its associated reward function, facilitating direct optimization for human preferences through a standard cross-entropy loss. This circumvents the need for reinforcement learning altogether and does so without sacrificing generality. Remarkably, DPO achieves results on par with, or even surpasses, those of prevalent RLHF techniques such as those employing PPO, while also demanding minimal hyperparameter adjustment. As a result, DPO substantially lowers the difficulty of preference-driven language model training.

\textbf{Limitations \& Future Work.} Our findings highlight several directions that warrant further investigation. One open question concerns the generalization capabilities of the DPO policy outside the training distribution, particularly in comparison to approaches that use explicit reward functions; preliminary findings indicate comparable generalization to PPO-based models, but more systematic exploration is required. Another aspect worth studying is whether leveraging self-generated labels from DPO-trained policies can similarly exploit unlabeled data efficiently. Furthermore, the implications of reward over-optimization within direct preference optimization remain unclear: for example, it is uncertain whether the modest performance dip observed in Figure~\ref{fig:dialogue-main}-right is a manifestation of such effects. Moreover, while our analysis includes models of up to 6B parameters, investigating the extension of DPO to much larger, state-of-the-art language models offers a promising avenue for research. On the evaluation front, we note that GPT-4’s computed win rates can depend on the prompt; future studies should examine optimal strategies for obtaining reliable judgments from automated evaluators. Ultimately, the potential applications of DPO extend beyond language model training, encompassing the broader training of generative models across multiple modalities.